% Separate post-CP21 candidate. All source/finite checks retained.
\subsection{The Newton filtration in the remaining reducible case}

Throughout B60--B61, $p\ge5$, $K$ is complete discretely valued of
mixed characteristic $(0,p)$ with algebraically closed residue field,
$e_K=p-1$, and $\mu_p\subset K$. We work in the normalized remaining
comparison of B49--B50: both groups are potentially supersingular of
height two, and their common finite determinant has been twisted to
the level-one tame fundamental character. Thus, putting
\[
              q=p^2,\qquad N=q-1,
\]
we may take the common good field $L/K$ to be tame cyclic of degree
$N$, and its absolute index is $e_L=N(p-1)$. Choose compatible tame
uniformizers and a primitive $N$th root $\zeta$ so that the generator
$\sigma$ acts on a uniformizer of $L$ by $\zeta$. An actual special
tangent exponent $c$ then satisfies
\[
                    c\equiv1\pmod{p-1}.                \tag{B-RD60.1}
\]
The special covariant character is the pair of exponents $c,pc$.
The normalization is a common finite height-one twist, as in B53;
it is not an identification of the mod-$p$ cyclotomic character
with the fundamental character.

\begin{lemma}[An elementary lower-break bound]
\label{lem:B-RD60-break}
\status{Candidate; complete proof written, independent audits pending.}
Let $E/D$ be a finite totally ramified Galois extension of complete
discretely valued fields with residue characteristic $p$, and let
$g$ be a nonidentity wild inertia element. Normalize $v_E$ to have
value group $\mathbf Z$, and put
\[
       b(g)=v_E(g\varpi-\varpi)-1
\]
for a uniformizer $\varpi$ of $E$. For every $x\in\mathfrak m_E$,
\[
 v_E(gx-x)\ge v_E(x)+b(g).                              \tag{B-RD60.2}
\]
Equality holds if $p\nmid v_E(x)$; if $p\mid v_E(x)$ the inequality
is strict. The statements allow $gx=x$, with valuation $+\infty$.
\end{lemma}
\begin{proof}
Choose representatives in $\mathcal O_D$ of every residue-field
element. Completeness gives an expansion
$x=\sum_{j\ge A}a_j\varpi^j$, with $A=v_E(x)$, all $a_j$ among
these fixed representatives and $a_A$ a unit. Write
$g\varpi=\varpi(1+d\varpi^{b(g)}+\cdots)$, with $d$ a unit.
The change of the degree-$j$ term has valuation at least $j+b(g)$.
If $p\nmid A$, its degree-$A$ term has exactly that valuation and
nonzero leading coefficient $Aa_A\bar d$; all higher terms change
later. If $p\mid A$, the linear binomial coefficient is in the
maximal ideal, while every higher binomial term has strictly larger
valuation because $b(g)\ge1$. The degree-$A$ change and all changes
from $j>A$ are therefore strictly later than $A+b(g)$. This proves
all assertions. No identification of lower and upper numbering is
used here.
\end{proof}

\begin{lemma}[The canonical high line and its two valuations]
\label{lem:B-RD60-Newton}
\status{Candidate; complete proof written, independent audits pending.}
Let $M$ be one normalized potentially supersingular group, and
suppose its residual inertia representation $M[p]$ is reducible over
$\mathbf F_p$. Its nonzero $p$-torsion points have exactly two
valuations in the good formal coordinate over $L$. There is a unique
high-valuation line $H\subset M[p]$ and an integer $B$ such that
\[
 \begin{gathered}
 0<B<p(p-1),\qquad
 v_L(T(Q))=\alpha=N-B\quad(0\ne Q\in H),\\
 v_L(T(Q))=\beta=B/p\quad(Q\in M[p]\setminus H),\qquad
 \alpha>\beta.                                        \tag{B-RD60.3}
 \end{gathered}
\]
The line $H$ is inertia-stable and is rational over $L$.
In the exact tame coordinate, its character on $\sigma$ is
\[
                  \zeta^{c-B},\qquad
                         B\equiv c\pmod{p+1}.         \tag{B-RD60.4}
\]
For every point $P$ of order $p^2$ with $pP\notin H$,
\[
                       v_L(T(P))=B/p^3.               \tag{B-RD60.5}
\]
All these valuations are unchanged by an integral invertible
formal-coordinate change.
\end{lemma}
\begin{proof}
Use the height-two universal coordinate and coefficient estimate
in Lemma~\ref{lem:B-RD49-point}, read there from the original
Hopkins--Gross Propositions~5.7 and~12.10
\cite[Proposition~5.7, equations (12.3)--(12.5), and
Proposition~12.10 with proof]{HopkinsGross1994Bundles}. Put $e=e_L$ and
$b=v_L(u)$, allowing $+\infty$. The only possible negative-slope
vertices for nonzero $p$-torsion are
\[
                        (1,e),\quad(p,b),\quad(q,0).
\]
Indeed coefficients in degrees $2,\ldots,p-1$ vanish in that
coordinate, coefficients strictly between $p$ and $q$ lie in
$(p,u)$, and the degree-$q$ coefficient is a unit.
If $b\ge ep/(p+1)$, all nonzero points have one valuation $e/N$.
Their leading-coordinate map is injective: equal leading
coefficients would give a nonzero difference of larger valuation,
again a $p$-torsion point. Wild inertia acts trivially on leading
coefficients after good reduction and hence on this entire residual
module. The remaining cyclic tame representation, acting as scalar
multiplication on its leading-coordinate plane in the residue field,
is either irreducible over $\mathbf F_p$ or scalar over
$\mathbf F_p$. The former contradicts the hypothesis; the latter
would give a square residual determinant. Our normalized determinant
is the nonsquare level-one character, so this case is impossible.

Consequently $0<b<ep/(p+1)$. The Newton polygon gives $p-1$ points
of valuation $(e-b)/(p-1)$ and $q-p$ points of valuation
$b/(p(p-1))$. The former, together with zero, form a line $H$:
formal addition and inversion preserve the condition of having at
least that valuation. Descent preserves valuations, by the same
integral good-model and Tate full-faithfulness argument as in B53,
so $H$ is inertia-stable. Its character has values in
$\mathbf F_p^\times$; it is tame of order dividing $p-1$ and is
therefore killed over $L$. Its coordinate valuation is an integer.
Writing that integer as $N-B$ gives $b=(p-1)B$, with $B$ an integer
in the range of (B-RD60.3).

For the character calculation use tame averaging, as in B56, to
choose an integral parameter in which the descent action on points
is $T\mapsto\zeta^c\sigma(T)$. In this parameter the coefficient
$a_p$ of degree $p$ in $[p](T)$ still has valuation $b<e$.
To check this last assertion without assuming a $p$-typical averaged
coordinate, reduce the coordinate change modulo $p$: the first
nonzero degree of $[p](T)$ is $p$, and its coefficient is changed
by a unit $(p-1)$st power. Terms from $pT$ lie in $p\mathcal O_L$.
Thus the valuation $b<e$ is unchanged. Compatibility with $[p]$
gives
\[
                  \sigma(a_p)=\zeta^{(p-1)c}a_p.
\]
Taking its leading coefficient in the tame uniformizer shows
$b\equiv(p-1)c\pmod N$, hence $B\equiv c\pmod{p+1}$.
For $0\ne Q\in H$, the coordinate valuation is $N-B$;
its field action contributes $\zeta^{N-B}$ and its descent action
contributes $\zeta^c$. This proves (B-RD60.4).

Finally $pP$ has valuation $\beta=B/p<p-1$. In the equation
$[p](T)=T(pP)$, its constant term and unit degree-$q$ term give a
single Newton segment: the linear term has coefficient valuation
$e>\beta$, the degree-$p$ coefficient has valuation
$b=(p-1)B>\beta$, and all intermediate coefficients are at least
as large as $b$. Every one of its $q$ roots thus has valuation
$\beta/q=B/p^3$. This is (B-RD60.5).
\end{proof}

\begin{proposition}[The first residual wild break]
\label{prop:B-RD60-residual-break}
\status{Candidate; complete proof written, independent audits pending.}
Let $R=L(M[p])$. Then $R=L$ or $R/L$ is cyclic of degree $p$.
In the former case $p\mid B$. In the latter case $H$ is the unique
line fixed by wild inertia and, writing $r$ for the single lower
break of $R/L$,
\[
 \begin{cases}
 r=pN-(p+1)B,&p\nmid B,\\
 r<pN-(p+1)B,&p\mid B.
 \end{cases}                                          \tag{B-RD60.6}
\]
If two normalized groups have isomorphic residual modules with
nontrivial wild image, then their high lines agree under that
isomorphism. If both associated integers $B_i$ are prime to $p$,
or both are divisible by $p$, their special characters already
agree. In the only remaining case, after indexing so that
$p\nmid B_1$ and $p\mid B_2$, one has $B_1>B_2$.
\end{proposition}
\begin{proof}
A reducible subgroup of $\GL_2(\mathbf F_p)$ is triangular. Its
wild image has order at most $p$, and its two tame diagonal
characters have order dividing $p-1$, killed over $L$. This proves
the assertion about $R$. If $R=L$, the low valuation $B/p$ is
integral, so $p\mid B$.

Otherwise the nontrivial unipotent group has one fixed line,
necessarily the invariant high line $H$. Every nonzero $g$ sends
a point $Q\notin H$ to $Q+Q_0$ for some $0\ne Q_0\in H$.
With $v_R$ normalized, its coordinate has valuation $B$, while
its difference has valuation $p(N-B)$. Ordinary coordinate
difference and formal-group subtraction have the same valuation:
the latter is the former times an integral unit. Applying
Lemma~\ref{lem:B-RD60-break} proves (B-RD60.6).

The common residual representation identifies the two unique high
lines and their characters. Thus
\[
                   c_1-B_1\equiv c_2-B_2\pmod N.       \tag{B-RD60.7}
\]
If both $B_i$ are prime to $p$, their common break gives $B_1=B_2$,
and (B-RD60.7) gives $c_1=c_2$ modulo $N$.
If both are divisible by $p$, write $B_i=pb_i$.
The range in (B-RD60.3) gives $1\le b_i\le p-2$.
Reducing (B-RD60.7) modulo $p-1$, and using (B-RD60.1), gives
$b_1\equiv b_2\pmod{p-1}$; the range forces equality, and again
the actual tangent exponents agree.
In the remaining case, equality of the common break and the strict
inequality on the second side give
\[
      pN-(p+1)B_1=r<pN-(p+1)B_2,
\]
so $B_1>B_2$ as asserted. No pair is constructed by these necessary
relations.
\end{proof}
